% =========================================================
% Computational Linear Algebra --- Independence, Basis, Dimension
% (Strang \S{}3.4 / Johnston \S{}2.3 / Ashraf \S{}2.2)
% =========================================================

\section{Independence, Basis, and Dimension}

% ---------------------------------------------------------
\subsection*{Linear Independence}

\begin{propositionbox}{Definition (Linear Independence)}
Vectors $\mathbf{v}_1, \mathbf{v}_2, \ldots, \mathbf{v}_k \in \mathbb{R}^n$ are
\emph{linearly independent} if the only solution to
\[
  c_1\mathbf{v}_1 + c_2\mathbf{v}_2 + \cdots + c_k\mathbf{v}_k = \mathbf{0}
\]
is $c_1 = c_2 = \cdots = c_k = 0$.
Otherwise they are \emph{linearly dependent}: at least one vector is
a linear combination of the others.
\end{propositionbox}

% ---------------------------------------------------------
\subsection*{Basis and Dimension}

\begin{propositionbox}{Definition (Basis)}
A \emph{basis} for a subspace $V \subseteq \mathbb{R}^n$ is a set of vectors
$\{\mathbf{v}_1, \ldots, \mathbf{v}_k\}$ that is:
\begin{enumerate}
  \item linearly independent, and
  \item spans $V$ (every vector in $V$ is a linear combination of $\mathbf{v}_1,\ldots,\mathbf{v}_k$).
\end{enumerate}
\end{propositionbox}

\begin{theorembox}{Theorem (Dimension is Well-Defined)}
Any two bases for the same subspace $V$ have the same number of vectors.
This common number is called the \emph{dimension} of $V$, written $\dim V$.
\end{theorembox}

\begin{remark*}
The standard basis $\{\mathbf{e}_1, \ldots, \mathbf{e}_n\}$ is a basis for $\mathbb{R}^n$,
so $\dim \mathbb{R}^n = n$.
The key insight: every set of $n$ linearly independent vectors in $\mathbb{R}^n$ is
automatically a basis.
\end{remark*}

% ---------------------------------------------------------
\subsection*{Change of Basis}

\begin{propositionbox}{Definition (Transition Matrix)}
Let $\mathcal{B}_1 = \{\mathbf{u}_1, \ldots, \mathbf{u}_n\}$ and
$\mathcal{B}_2 = \{\mathbf{v}_1, \ldots, \mathbf{v}_n\}$ be two ordered bases for $\mathbb{R}^n$.
The \emph{transition matrix} $P$ from $\mathcal{B}_1$ to $\mathcal{B}_2$ has
as its $i$-th column the coordinates of $\mathbf{u}_i$ expressed in the basis $\mathcal{B}_2$.
If $[\mathbf{x}]_{\mathcal{B}_1}$ denotes the coordinate vector of $\mathbf{x}$ in basis $\mathcal{B}_1$, then
\[
  [\mathbf{x}]_{\mathcal{B}_2} = P \,[\mathbf{x}]_{\mathcal{B}_1}.
\]
The inverse $P^{-1}$ converts in the opposite direction.
\end{propositionbox}

\begin{remark*}[NURBS connection: Bernstein basis as a basis for polynomials]
The Bernstein polynomials $\{B_{0,n}(u), B_{1,n}(u), \ldots, B_{n,n}(u)\}$ form a
basis for the vector space of polynomials of degree $\leq n$.
So does the monomial (power) basis $\{1, u, u^2, \ldots, u^n\}$.
The matrix $M$ from Piegl \& Tiller Exercise 1.14 is precisely the transition
matrix converting coordinates in the monomial basis to coordinates in the Bernstein basis.
$M^{-1}$ converts the other direction.
This is why the B\'{e}zier curve $C(u) = [u^i]^T M [P_i]$ can be read as:
``express the Bernstein basis functions in the monomial basis via $M$, then
multiply by the control points.''
\end{remark*}
